\section{Attainment under every full-support prior}
\label{sec:v18-universal}
The rank formula uses the specified prior. We now characterize when a
square product pairing has full normalized rank for \emph{every}
full-support prior. This replaces separate Chebyshev assumptions by an
exact sign condition on the paired evaluation determinants. The condition
is intrinsic to the two function spaces and does not require an ordering
of $X$.

\subsection{Positive histories and their augmented tangents}
Fix an interior command tuple $u_*$ for a positive history product
$P(u,\cdot)\in C(X)$, twice continuously differentiable in $u$. Write
$P=P(u_*,\cdot)>0$, and let
\begin{equation}\label{eq:v18-augmented-tangent}
 \widetilde T=\operatorname{span}\{P\}+\operatorname{im}DP(u_*).
\end{equation}
Let $p\ge1$ and $\Phi=\operatorname{span}\{1,\phi_1,\ldots,\phi_p\}$, where the
functions are real and continuous, and suppose
\begin{equation}\label{eq:v18-square-pair}
 \dim\widetilde T=\dim\Phi=p+1=:d.
\end{equation}
This is the square pairing regime; no claim that every rectangular
pairing reduces to it is made. For a basis $v_0,\ldots,v_p$ of
$\widetilde T$, with $\phi_0=1$, define the continuous symmetric
function on $X^d$
\begin{equation}\label{eq:v18-sign-kernel}
 \Delta_{v,\phi}(x_0,\ldots,x_p)
 =\det[v_i(x_k)]_{i,k=0}^p\,
  \det[\phi_j(x_k)]_{j,k=0}^p.
\end{equation}
It is symmetric because a permutation changes both determinants by the
same sign. Changing either basis multiplies it by a nonzero constant;
being of one weak sign and nonzero somewhere is therefore a property of
the pair $(\widetilde T,\Phi)$, not of its oriented bases.

\begin{theorem}[Universal square-pairing criterion]
\label{thm:v18-universal-pairing}
Under \eqref{eq:v18-square-pair}, the following assertions are equivalent.

\smallskip
\noindent\textup{(i)} For every full-support probability $\mu$ on $X$,
the actual command map
\[
 u\longmapsto\bigl(\mu(P(u)\phi_j)/\mu P(u)\bigr)_{j=1}^p
\]
has rank $p$ at $u_*$.

\smallskip
\noindent\textup{(ii)} The function $\Delta_{v,\phi}$ has one weak
sign on $X^d$ and is nonzero at at least one tuple.

\smallskip
Under these equivalent conditions, independent commands with a density
bounded below near $u_*$ give a positive subprobability minorization of
a $p$-dimensional prediction cube for every such prior. If the positive
history products, tests, interior tuples, and command derivatives vary
continuously in a compact family, the dimensions in
\eqref{eq:v18-square-pair} remain fixed, and \textup{(ii)} holds at
every parameter, the minorization is uniform over that family and all
probabilities $\mu\ge c\mu_0$, where $c>0$ and $\mu_0$ has full support.
A common positive lower bound for the likelihoods, a common command
neighborhood, and a common lower command-density bound are understood
in this uniform assertion.
\end{theorem}
\begin{proof}
For each $\mu$ define $L_\mu Q=(\mu Q,\mu(Q\phi_1),\ldots,
\mu(Q\phi_p))^t$ and $Z=\mu P$. The normalized differential is the
map
\[
 Q\longmapsto Z^{-1}\left(L_\mu Q-
              \frac{L_\mu P}{Z}\,\mu Q\right)
\]
on actual product derivatives, with its constant coordinate omitted.
Adding the direction $P$ does not change its image, because that
direction maps to zero. On $L_\mu(\widetilde T)$ the kernel of the
projection in parentheses is exactly the span of $L_\mu P$, whose
constant coordinate is $Z>0$. Therefore the normalized command rank
is $\rank(L_\mu|\widetilde T)-1$. In the square case it equals $p$
if and only if
\[
 D(\mu):=\det[\mu(v_i\phi_j)]_{i,j=0}^p\ne0.
\]
This argument explains why adjoining the evidence direction in
\eqref{eq:v18-augmented-tangent} is legitimate even when $P$ is not an
actual command derivative.

Determinant expansion, Fubini, and symmetrization give the classical
integration identity
\begin{equation}\label{eq:v18-andreief}
 D(\mu)=\frac1{d!}\int_{X^d}\Delta_{v,\phi}(x_0,\ldots,x_p)
                     \prod_{k=0}^p\mu(dx_k).
\end{equation}
For completeness, expand both determinants on the right into their
permutation sums. Integrating a term gives a product of $d$ scalar
integrals. Reindexing one permutation relative to the other produces
the determinant on the left exactly $d!$ times. All functions are
bounded, so the expansions and integrations are legitimate. This is
Andr\'eief's identity, not a new integration formula; see
\cite{ForresterAndreief} for its source and attribution.

Suppose \textup{(ii)} holds and orient a basis so that $\Delta\ge0$.
Continuity and a strictly positive value give a product of nonempty
open sets on which $\Delta>0$. Every full-support probability gives
this product positive mass. Equation~\eqref{eq:v18-andreief} therefore
gives $D(\mu)>0$, proving \textup{(i)}.

Conversely, the full-support probabilities form a convex set. If
\textup{(i)} holds, continuity along every line segment of this set
shows that $D(\mu)$ has the same nonzero sign throughout it. Orient a
basis so that the sign is positive and fix one full-support $\mu_0$.
For an arbitrary tuple $(x_0,\ldots,x_p)$ put
$\nu=d^{-1}\sum_k\delta_{x_k}$. The probabilities
$(1-\epsilon)\nu+\epsilon\mu_0$ have full support for every
$0<\epsilon<1$. Letting $\epsilon\downarrow0$ gives $D(\nu)\ge0$.
Matrix multiplication at this atomic probability gives, also for
repeated points,
\[
 D(\nu)=d^{-d}\det[v_i(x_k)]\det[\phi_j(x_k)]
       =d^{-d}\Delta_{v,\phi}(x_0,\ldots,x_p).
\]
Hence $\Delta\ge0$ everywhere. Since $D(\mu_0)>0$ in
\eqref{eq:v18-andreief}, $\Delta$ is not identically zero. This proves
\textup{(ii)} and its necessity without assuming either evaluation
determinant separately has a fixed sign.

The minorization follows by augmenting the onto normalized derivative
with orthonormal kernel coordinates, applying the inverse-function
theorem, and integrating the entire kernel-coordinate cube, exactly as
in the proof of Lemma~\ref{lem:v18-image-mass}. The density of the
specified report word contains its evidence, bounded below by the
product of the likelihood lower bounds.

For uniformity, the set of probabilities dominating $c\mu_0$ is weakly
compact: it is a closed subset of the probabilities on compact $X$,
as is seen by testing the domination inequality against continuous
nonnegative functions. Every limiting probability still has full
support. The normalized command Jacobian is jointly continuous in the
parameter and this weakly varying probability; this follows from the
uniform-norm continuity of the finitely many integrands and the positive
evidence bound. It is onto everywhere by the equivalence just proved.
Compactness gives a positive lower bound for its least row singular
value. Uniform second derivative bounds and uniformly interior tuples
then give uniform inverse-chart radii, while the specified command and
evidence lower bounds give uniform minorization mass. No conclusion
about a later observation operator's small singular values is used.
\end{proof}

\begin{corollary}[Matched score and query spaces on arbitrary latent spaces]
\label{cor:v18-matched-spaces}
If $\widetilde T=P\Phi$ and $\dim\Phi=p+1$, then the normalized
history map has rank $p$ for every full-support prior. More generally,
the same conclusion holds whenever the paired evaluation determinants
have one weak sign and are nonzero somewhere, even if both individual
determinants change sign.
\end{corollary}
\begin{proof}
Choose $v_i=P\phi_i$. Then
\[
 \Delta_{v,\phi}(x_0,\ldots,x_p)
   =\left(\prod_{k=0}^pP(x_k)\right)
     \det[\phi_i(x_k)]_{i,k=0}^p{}^2\ge0.
\]
The evaluation determinant is nonzero somewhere. Indeed, if the
vectors $(\phi_i(x))_{i=0}^p$ did not span $\R^{p+1}$ as $x$ varies,
a nonzero linear functional annihilating all of them would contradict
the independence of the functions. Selecting a spanning set of $p+1$
evaluation vectors gives the desired tuple. Apply
Theorem~\ref{thm:v18-universal-pairing}. Its same equivalence proves
the more general assertion.
\end{proof}

\subsection{Relation to positive products and to rank strata}
For the positive products of Section~\ref{sec:positive-history}, the
original normalized-product criterion remains valid. The present
criterion makes two further distinctions. It uses the actual augmented
tangent, so it does not require a physically available radial command.
It also characterizes, in the square regime, the full-support-prior
quantifier itself. Two separately strict Chebyshev systems with compatible
orientations make \eqref{eq:v18-sign-kernel} positive on distinct ordered
tuples and are a sufficient instance. They are not necessary: the
matched-space corollary works on compact spaces of any dimension.

For example, at $u=0$ in \eqref{eq:v18-affine-likelihood}, the augmented
tangent is $\operatorname{span}\{1,f_1,\ldots,f_b\}$. When the future
tests span that same space and these functions are independent, the
corollary proves maximal acquisition for every full-support prior.
In contrast, for $v=(1,x)$ and $\phi=(1,x^2)$ on $[-1,1]$, the paired
evaluation determinant is
\[
 (x_1-x_0)^2(x_0+x_1),
\]
which takes both signs. The normalized rank can consequently depend on
the full-support prior. Proposition~\ref{prop:v18-interior-example}
computes that dependence and its entire finite-label profile exactly.
These observations link the universal criterion, the algebraic moment
strata, and the covariance classification without substituting a local
pairing argument for either a global cover or a causal construction.
